
\subsection{Caso de estudio: una sala, paso a paso}\label{sec:sala}

Igual que con el casco, el proceso de las habitaciones se desarrolla aquí con una fotografía real y con los números que produce el programa. La
foto es una sala con un sofá, una mesa de centro y un mueble de televisión (figura~\ref{fig:hab-entrada}), tomada de Wikimedia Commons
(autor: Shixart1985, licencia CC BY 2.0, \texttt{commons.wikimedia.org}). Se reduce a $960\times640$ píxeles ($N_p=614\,400$ píxeles) y se procesa con el
programa tal como está (procesamiento bajo); cada etapa se ejecutó por separado
(\texttt{experimentos/habitacion\_paso\_a\_paso.py}). En la figura~\ref{fig:hab-flujo} están los seis pasos y los objetos matemáticos de cada uno.

\begin{figure}[htbp]
\centering
\includegraphics[width=0.95\linewidth]{hab_entrada.jpg}
\caption{Foto de la sala reducida a $960\times640$ píxeles, con los ejes $u$ (columna) y $v$ (fila), el centro $(c_x,c_y)$ y los cuatro píxeles
(piso, pared, techo y sofá) que se siguen en las cuentas.}
\label{fig:hab-entrada}
\end{figure}

\begin{figure}[htbp]
\centering
\includegraphics[width=\linewidth]{hab_flujo.jpg}
\caption{Los seis pasos de la reconstrucción de una habitación y el objeto matemático de cada uno.}
\label{fig:hab-flujo}
\end{figure}

\paragraph{La cámara de la sala.} Con el ángulo de visión $\alpha=65^\circ$ y $W=960$, la distancia focal es $f=(W/2)/\tan(\alpha/2)=753{,}449$ píxeles,
$(c_x,c_y)=(480,320)$ y
\begin{equation}
K=\begin{pmatrix} 753{,}449 & 0 & 480 \\ 0 & 753{,}449 & 320 \\ 0 & 0 & 1 \end{pmatrix},\qquad \det K=f^2=567\,685{,}5,
\end{equation}
\begin{equation}
K^{-1}=\begin{pmatrix} 0{,}001327 & 0 & -0{,}637070 \\ 0 & 0{,}001327 & -0{,}424714 \\ 0 & 0 & 1 \end{pmatrix} .
\end{equation}
Son los mismos valores que en el casco porque la foto tiene el mismo ancho; solo cambia $c_y$ porque la altura es otra.

\paragraph{Paso 1: profundidad en metros y clase de cada píxel.} La red Depth Anything V2 entrega una matriz $Z$ de $640\times960$ números con
la distancia en metros de cada píxel al plano de la cámara: va de 1{,}40 m (la mesa de centro) a 7{,}50 m (el fondo), con media de 4{,}38 m,
y tardó 1\,583 ms. SegFormer entrega otra matriz del mismo tamaño con la clase de cada píxel (24 clases distintas en esta foto, en 3\,692 ms); las
del cuarto se cuentan en el cuadro~\ref{tab:hab-clases} y las dos matrices se ven en la figura~\ref{fig:hab-etapas}.

\begin{table}[htbp]
\caption{Píxeles de la foto de la sala que pertenecen a cada superficie del cuarto}
\label{tab:hab-clases}
\centering\small
\begin{tabular}{@{}lrr@{}}
\toprule
Clase & Píxeles & Parte de la foto \\
\midrule
Pared & $186\,555$ & $30{,}4$\,\% \\
Piso & $61\,644$ & $10{,}0$\,\% \\
Techo & $47\,065$ & $7{,}7$\,\% \\
Ventana & $36\,010$ & $5{,}9$\,\% \\
Espejo & $8\,894$ & $1{,}4$\,\% \\
\bottomrule
\end{tabular}
\end{table}

\begin{figure}[htbp]
\centering
\includegraphics[width=0.97\linewidth]{hab_etapas.jpg}
\caption{De la foto a las matrices del proceso: foto, profundidad $Z(u,v)$ de la red, clase de cada píxel y planos ajustados (un color por plano).}
\label{fig:hab-etapas}
\end{figure}

\paragraph{De los píxeles a vectores.} Cada píxel se convierte en un vector con $\vect{r}=K^{-1}\vect{p}$ (como en la
sección «\nameref{sec:vectores-casco}») y luego en un punto $\vect{P}=Z\,\vect{r}$. El cuadro~\ref{tab:hab-rayos} muestra los cuatro píxeles de la figura~\ref{fig:hab-entrada},
y la retroproyección de los cuatro a la vez es un solo producto de matrices, $P=K^{-1}H\operatorname{diag}(Z_1,\dots,Z_4)$, con $H$ de $3\times4$ (cada
columna es un píxel):
\begin{equation}
H=\begin{pmatrix} 631 & 154 & 432 & 173 \\ 522 & 288 & 26 & 461 \\ 1 & 1 & 1 & 1 \end{pmatrix},
\end{equation}
\begin{equation}
K^{-1}H=\begin{pmatrix} 0{,}2003 & -0{,}4327 & -0{,}0638 & -0{,}4075 \\ 0{,}2680 & -0{,}0425 & -0{,}3902 & 0{,}1870 \\ 1 & 1 & 1 & 1 \end{pmatrix},
\end{equation}
\begin{equation}
P=K^{-1}H\begin{pmatrix}3{,}85&0&0&0\\0&4{,}70&0&0\\0&0&5{,}02&0\\0&0&0&2{,}79\end{pmatrix}=\begin{pmatrix} 0{,}771 & -2{,}033 & -0{,}321 & -1{,}136 \\ 1{,}032 & -0{,}200 & -1{,}960 & 0{,}521 \\ 3{,}851 & 4{,}698 & 5{,}023 & 2{,}788 \end{pmatrix}\ \text{m} .
\end{equation}
Las columnas de $P$ son los puntos del piso, de la pared, del techo y del sofá en los ejes de la cámara.

\begin{table}[htbp]
\caption{Cuatro píxeles de la sala como vectores: píxel $(u,v)$, rayo $\vect{r}=K^{-1}\vect{p}$, su norma, el ángulo con el eje óptico y la profundidad $Z$ de la red (m)}
\label{tab:hab-rayos}
\centering\small
\begin{tabular}{@{}llcccc@{}}
\toprule
 & $(u,v)$ & $\vect{r}=K^{-1}\vect{p}$ & $\lVert\vect{r}\rVert$ & $\theta$ & $Z$ \\
\midrule
Piso & $(631,\,522)$ & $(0{,}2004,\,0{,}2681,\,1)$ & $1{,}055$ & $18{,}5^\circ$ & $3{,}85$ \\
Pared & $(154,\,288)$ & $(-0{,}4327,\,-0{,}0425,\,1)$ & $1{,}090$ & $23{,}5^\circ$ & $4{,}70$ \\
Techo & $(432,\,26)$ & $(-0{,}0637,\,-0{,}3902,\,1)$ & $1{,}075$ & $21{,}6^\circ$ & $5{,}02$ \\
Sofá & $(173,\,461)$ & $(-0{,}4075,\,0{,}1871,\,1)$ & $1{,}096$ & $24{,}2^\circ$ & $2{,}79$ \\
\bottomrule
\end{tabular}
\end{table}

\paragraph{Paso 2: ajustar planos a los puntos con RANSAC.} El piso, el techo y las paredes son planos, es decir, conjuntos de puntos $\vect{P}$ que cumplen
$\vect{n}\cdot\vect{P}=d$ con $\vect{n}$ de norma $1$. Como la red tiene ruido, los puntos no caen exactamente en un plano, y se busca el
plano que más puntos explica con el método RANSAC, que repite tres pasos:
(1) tomar tres puntos al azar de la misma clase, (2) calcular el plano que pasa por ellos con el producto cruz y (3) contar cuántos de los demás puntos
están a menos de una tolerancia de ese plano. Para el piso, el primer intento con una muestra de $20\,000$ puntos de la clase «piso» tomó los
píxeles $(456,\,423)$, $(687,\,538)$ y $(501,\,415)$, que en el espacio son
\begin{equation}
\vect{P}_0=\begin{pmatrix} -0{,}199 \\ 0{,}853 \\ 6{,}243 \end{pmatrix},\quad \vect{P}_1=\begin{pmatrix} 0{,}952 \\ 1{,}003 \\ 3{,}465 \end{pmatrix},\quad \vect{P}_2=\begin{pmatrix} 0{,}182 \\ 0{,}823 \\ 6{,}527 \end{pmatrix}\ \text{m},
\end{equation}
y dos lados del triángulo que forman son
\begin{equation}
\vect{e}_1=\vect{P}_1-\vect{P}_0=\begin{pmatrix} 1{,}151 \\ 0{,}149 \\ -2{,}778 \end{pmatrix},\qquad \vect{e}_2=\vect{P}_2-\vect{P}_0=\begin{pmatrix} 0{,}381 \\ -0{,}030 \\ 0{,}284 \end{pmatrix} .
\end{equation}
El producto cruz $\vect{n}=\vect{e}_1\times\vect{e}_2$ es perpendicular a los dos y, por tanto, al plano. Sus componentes son
\begin{align*}
n_x&=e_{1y}e_{2z}-e_{1z}e_{2y}=0{,}149\cdot0{,}284-\left(-2{,}778\right)\cdot\left(-0{,}030\right)=-0{,}0410,\\
n_y&=e_{1z}e_{2x}-e_{1x}e_{2z}=\left(-2{,}778\right)\cdot0{,}381-1{,}151\cdot0{,}284=-1{,}3853,\\
n_z&=e_{1x}e_{2y}-e_{1y}e_{2x}=1{,}151\cdot\left(-0{,}030\right)-0{,}149\cdot0{,}381=-0{,}0913,
\end{align*}
es decir $\vect{n}=\begin{pmatrix} -0{,}041 \\ -1{,}385 \\ -0{,}091 \end{pmatrix}$, con norma $\lVert\vect{n}\rVert=1{,}389$. Se divide entre la norma para que quede de longitud $1$ y se elige el signo
para que apunte hacia arriba (el eje $y$ de la cámara apunta hacia abajo), y $d$ sale del producto punto con cualquiera de los tres puntos:
\begin{equation}
\hat{\vect{n}}=\frac{\vect{n}}{\lVert\vect{n}\rVert}=\begin{pmatrix} -0{,}0295 \\ -0{,}9974 \\ -0{,}0657 \end{pmatrix},\qquad d=\hat{\vect{n}}\cdot\vect{P}_0=-1{,}255\ \text{m}.
\end{equation}
La tolerancia es el 2\,\% de la mediana de $Z$ de la muestra, 7{,}0 cm, y 19\,413 de los 20\,000 puntos de la muestra quedaron dentro de ella
(97{,}1\,\%): el plano es muy bueno. Con todos esos puntos se vuelve a ajustar el plano por mínimos cuadrados (el vector singular del menor valor
singular de los puntos centrados, el mismo SVD que se usa en la sección «\nameref{sec:inversa}») y queda el plano final del piso. El método repite esto
para el techo y hasta cuatro paredes. El cuadro~\ref{tab:hab-planos} reúne los seis planos y las figuras~\ref{fig:hab-ransac} y~\ref{fig:hab-planos} los muestran.

\begin{figure}[htbp]
\centering
\includegraphics[width=0.97\linewidth]{hab_ransac.jpg}
\caption{Izquierda: tres puntos del piso, los lados $\vect{e}_1$ y $\vect{e}_2$ de su triángulo y la normal $\hat{\vect{n}}$ del producto cruz.
Derecha: distancia de los puntos de la muestra a ese plano; los que quedan a la izquierda de la línea roja votan por él.}
\label{fig:hab-ransac}
\end{figure}

\begin{table}[htbp]
\caption{Planos ajustados en la sala: normal unitaria $\hat{\vect{n}}$, constante $d$ (m), ángulo de la normal con la vertical y píxeles que quedaron sobre el plano}
\label{tab:hab-planos}
\centering\small
\begin{tabular}{@{}lp{5.2cm}rrr@{}}
\toprule
Plano & $\hat{\vect{n}}=(n_x,n_y,n_z)$ & $d$ & Ángulo & Píxeles \\
\midrule
Piso & $(0{,}0100,\,-0{,}9982,\,-0{,}0587)$ & $-1{,}210$ & $3{,}4^\circ$ & $61\,102$ \\
Techo & $(-0{,}0073,\,-0{,}9997,\,0{,}0217)$ & $2{,}088$ & $1{,}3^\circ$ & $47\,065$ \\
Pared 1 & $(-0{,}9993,\,-0{,}0184,\,0{,}0309)$ & $2{,}191$ & $88{,}9^\circ$ & $84\,756$ \\
Pared 2 & $(-0{,}9937,\,-0{,}1087,\,-0{,}0263)$ & $-1{,}788$ & $83{,}8^\circ$ & $73\,902$ \\
Pared 3 & $(-0{,}2504,\,-0{,}1255,\,0{,}9600)$ & $7{,}079$ & $82{,}8^\circ$ & $50\,474$ \\
Pared 4 & $(-0{,}9937,\,-0{,}0270,\,0{,}1089)$ & $-1{,}240$ & $88{,}5^\circ$ & $12\,566$ \\
\bottomrule
\end{tabular}
\end{table}

\begin{figure}[htbp]
\centering
\includegraphics[width=\linewidth]{hab_planos.jpg}
\caption{Los puntos de cada plano, vistos desde arriba y de lado, con la cámara en el origen. El piso (azul) sube suavemente hacia el fondo, el techo
(naranja) y las paredes forman las otras caras del cuarto.}
\label{fig:hab-planos}
\end{figure}

\paragraph{Qué dicen las normales.} Las normales son vectores y el producto punto dice cómo se cruzan los planos. La del piso es
$\hat{\vect{n}}_{\text{piso}}=\begin{pmatrix} 0{,}0100 \\ -0{,}9982 \\ -0{,}0587 \end{pmatrix}$; la de la pared de la izquierda, $\hat{\vect{n}}_1=\begin{pmatrix} -0{,}9993 \\ -0{,}0184 \\ 0{,}0309 \end{pmatrix}$. Su producto punto es
\begin{equation}
\hat{\vect{n}}_{\text{piso}}\cdot\hat{\vect{n}}_1=0{,}0066\quad\Longrightarrow\quad\theta=\arccos(0{,}0066)=89{,}6^\circ,
\end{equation}
casi perpendiculares, como deben ser un piso y una pared. Entre la pared izquierda y la del fondo, $\hat{\vect{n}}_1\cdot\hat{\vect{n}}_3=0{,}282$ y el ángulo es
$73{,}6^\circ$: no son $90^\circ$ exactos porque el programa supone un ángulo de visión de $65^\circ$ y la red tiene error, y eso inclina un poco los
planos lejanos (figura~\ref{fig:hab-normales}). El techo es casi paralelo al piso ($4{,}7^\circ$). Además, el valor $|d|$ de la normal del piso es la
\emph{altura de la cámara}: $|d|=1{,}210$ m.

\begin{figure}[htbp]
\centering
\includegraphics[width=0.97\linewidth]{hab_normales.jpg}
\caption{Las normales de cuatro planos como vectores unitarios (izquierda) y el ángulo entre cada par de planos, calculado con el producto punto (derecha).}
\label{fig:hab-normales}
\end{figure}

\paragraph{Paso 3: intersección rayo--plano.} Para que los planos queden perfectamente planos, el píxel cambia su profundidad por la del plano. El
punto del plano sobre el rayo es $\vect{P}=t\,\vect{r}$ con $\hat{\vect{n}}\cdot\vect{P}=d$, y de ahí $t\,(\hat{\vect{n}}\cdot\vect{r})=d$, es decir
\begin{equation}
Z=\frac{d}{\hat{\vect{n}}\cdot\vect{r}} .
\end{equation}
Para el píxel del piso, $\vect{r}=\begin{pmatrix} 0{,}2004 \\ 0{,}2681 \\ 1 \end{pmatrix}$ y $\hat{\vect{n}}\cdot\vect{r}=(0{,}0100)(0{,}2004)+(-0{,}9982)(0{,}2681)+(-0{,}0587)(1)=-0{,}3244$, así que
\begin{equation}
Z=\frac{-1{,}210}{-0{,}3244}=3{,}730\ \text{m},\qquad \vect{P}=3{,}730\cdot\begin{pmatrix} 0{,}2004 \\ 0{,}2681 \\ 1 \end{pmatrix}=\begin{pmatrix} 0{,}748 \\ 1 \\ 3{,}730 \end{pmatrix}\ \text{m},
\end{equation}
frente a $Z=3{,}851$ m de la red (una diferencia de 3{,}1\,\%). La figura~\ref{fig:hab-rayoplano} lo muestra en un corte lateral. El cuadro~\ref{tab:hab-int}
repite la cuenta para el piso, la pared y el techo. Las cuentas solo valen si $\hat{\vect{n}}\cdot\vect{r}\neq0$ (el rayo no es paralelo al plano) y $Z>0$ (el
plano está delante de la cámara).

\begin{figure}[htbp]
\centering
\includegraphics[width=0.95\linewidth]{hab_rayo_plano.jpg}
\caption{Corte lateral de la sala: el rayo $\vect{r}=K^{-1}\vect{p}$ del píxel del piso corta el plano $\hat{\vect{n}}\cdot\vect{P}=d$ en
$Z=d/(\hat{\vect{n}}\cdot\vect{r})$; la distancia $|d|$ del origen al plano es la altura de la cámara.}
\label{fig:hab-rayoplano}
\end{figure}

\begin{table}[htbp]
\caption{Intersección rayo--plano para tres píxeles: $\hat{\vect{n}}\cdot\vect{r}$, $d$ (m), $Z=d/(\hat{\vect{n}}\cdot\vect{r})$ (m), $Z$ de la red (m) y distancia entre el punto de la red y el del plano (cm)}
\label{tab:hab-int}
\centering\small
\begin{tabular}{@{}llrrrrr@{}}
\toprule
 & Plano & $\hat{\vect{n}}\cdot\vect{r}$ & $d$ & $Z$ del plano & $Z$ de la red & Dist. \\
\midrule
Piso & piso & $-0{,}3244$ & $-1{,}210$ & $3{,}730$ & $3{,}851$ & $12{,}7$ \\
Pared & pared 1 & $0{,}4641$ & $2{,}191$ & $4{,}721$ & $4{,}698$ & $2{,}5$ \\
Techo & techo & $0{,}4123$ & $2{,}088$ & $5{,}064$ & $5{,}023$ & $4{,}3$ \\
\bottomrule
\end{tabular}
\end{table}

\paragraph{La misma cuenta como sistema de ecuaciones.} Hallar $\vect{P}=(X,Y,Z)$ en el rayo del píxel $(u,v)$ y en el plano $\hat{\vect{n}}\cdot\vect{P}=d$ es
un sistema lineal de $3\times3$: las dos filas de arriba son la proyección $K\vect{P}=Z\vect{p}$ reescrita, $fX+(c_x-u)Z=0$ y $fY+(c_y-v)Z=0$, y la de abajo es el plano.
Para el píxel del piso,
\begin{equation}
\underbrace{\begin{pmatrix}753{,}449&0&-151\\0&753{,}449&-202\\0{,}0100&-0{,}9982&-0{,}0587\end{pmatrix}}_{A}
\begin{pmatrix}X\\Y\\Z\end{pmatrix}=\begin{pmatrix}0\\0\\-1{,}210\end{pmatrix}.
\end{equation}
Su determinante es $\det A=f^{2}\,(\hat{\vect{n}}\cdot\vect{r})=(567\,685)(-0{,}3244)=-184\,145{,}3$ (se desarrolla por la primera columna y sale $f^2$ por el producto punto), que es distinto de
cero justamente cuando el rayo no es paralelo al plano; entonces $A$ es invertible y $\vect{P}=A^{-1}\vect{b}=\begin{pmatrix} 0{,}748 \\ 1{,}000 \\ 3{,}730 \end{pmatrix}$ m, el mismo punto de arriba. Es decir: la intersección
rayo--plano es un sistema $3\times3$ con solución única, y la fórmula $Z=d/(\hat{\vect{n}}\cdot\vect{r})$ es su solución despejada.

\begin{figure}[htbp]
\centering
\includegraphics[width=0.97\linewidth]{hab_planitud.jpg}
\caption{Izquierda: desviación de los puntos de cada superficie respecto de su plano antes (puntos de la red) y después de la intersección rayo--plano (escala
logarítmica). Derecha: cuánto cambia la profundidad, en promedio, al reemplazarla por la del plano.}
\label{fig:hab-planitud}
\end{figure}

La figura~\ref{fig:hab-planitud} resume el efecto en los seis planos: antes, los puntos del piso se desviaban 0{,}9 cm de su plano y los de las paredes entre 2{,}5 y 14{,}8 cm; después, la
desviación es menor que $10^{-3}$ cm (cero, hasta el redondeo). La profundidad de la red cambia entre 0{,}8\,\% y 2{,}8\,\% al pasar al plano.

\paragraph{El cuarto vacío.} Con los mismos planos se calcula qué hay \emph{detrás} de cada mueble. Para un píxel tapado, el rayo corta todos los planos del cuarto en
$t_k=d_k/(\hat{\vect{n}}_k\cdot\vect{r})$, y los valores negativos (planos detrás de la cámara) se descartan; el verdadero fondo es el menor $t_k$ positivo. Para el píxel del sofá,
$\vect{r}=\begin{pmatrix} -0{,}4075 \\ 0{,}1871 \\ 1 \end{pmatrix}$ y el cuadro~\ref{tab:hab-vacio} da los seis cortes; el menor positivo es 4{,}85 m (el piso), muy lejos del sofá ($Z=2{,}79$ m),
de modo que detrás del sofá se pone el piso a esa profundidad (figura~\ref{fig:hab-vacio}).

\begin{table}[htbp]
\caption{Cortes del rayo del píxel del sofá con cada plano: $t_k=d_k/(\hat{\vect{n}}_k\cdot\vect{r})$ y si queda delante de la cámara}
\label{tab:hab-vacio}
\centering\small
\begin{tabular}{@{}lrrrr@{}}
\toprule
Plano & $\hat{\vect{n}}_k\cdot\vect{r}$ & $d_k$ (m) & $t_k$ (m) & $t_k>0$ \\
\midrule
Piso & $-0{,}250$ & $-1{,}210$ & $4{,}85$ & sí \\
Techo & $-0{,}162$ & $2{,}088$ & $-12{,}85$ & no \\
Pared 1 & $0{,}435$ & $2{,}191$ & $5{,}04$ & sí \\
Pared 2 & $0{,}358$ & $-1{,}788$ & $-4{,}99$ & no \\
Pared 3 & $1{,}039$ & $7{,}079$ & $6{,}82$ & sí \\
Pared 4 & $0{,}509$ & $-1{,}240$ & $-2{,}44$ & no \\
\bottomrule
\end{tabular}
\end{table}

\begin{figure}[htbp]
\centering
\includegraphics[width=0.97\linewidth]{hab_vacio.jpg}
\caption{Izquierda: cortes del rayo del píxel del sofá con los seis planos; se elige el menor positivo. Derecha: vista desde arriba del rayo, el sofá y el
fondo que queda detrás de él.}
\label{fig:hab-vacio}
\end{figure}

\paragraph{Paso 4: la malla y los triángulos que se cortan.} Sobre la foto se coloca una rejilla de un nodo cada 2 píxeles: $480\times320=153\,600$ nodos que,
llevados al espacio con $\vect{P}=Z\,K^{-1}\vect{p}$, son las filas de la matriz de vértices $V$. Cada celda se divide en dos triángulos y en total hay
$2(480-1)(320-1)=305\,602$ triángulos. Donde un mueble tapa una pared, $Z$ salta de golpe (figura~\ref{fig:hab-corte}, izquierda) y los triángulos que unen los dos lados
quedan estirados, como cortinas, y vistos casi de canto. Un triángulo se elimina si cumple las dos condiciones siguientes, calculadas con vectores:
\begin{itemize}
  \item se ve de canto: $\lvert\hat{\vect{n}}\cdot\hat{\vect{c}}_t\rvert<\operatorname{sen}5^\circ=0{,}087$, con $\hat{\vect{n}}$ su normal (producto cruz de dos lados) y $\hat{\vect{c}}_t$ el
        rayo hacia su centro, y
  \item la profundidad salta: $Z_{\max}/Z_{\min}>1{,}12$.
\end{itemize}
Con un triángulo del piso, de vértices en los píxeles $(576,548)$, $(576,550)$, $(578,548)$, los lados (en mm) son $\vect{e}_1=\begin{pmatrix} -2 \\ 1{,}1 \\ -16 \end{pmatrix}$ y $\vect{e}_2=\begin{pmatrix} 6{,}7 \\ 1{,}7 \\ 5{,}5 \end{pmatrix}$ y
\begin{align*}
n_x&=e_{1y}e_{2z}-e_{1z}e_{2y}=1{,}1\cdot5{,}5-\left(-16{,}0\right)\cdot1{,}7=33{,}25,\\
n_y&=e_{1z}e_{2x}-e_{1x}e_{2z}=\left(-16{,}0\right)\cdot6{,}7-\left(-2{,}0\right)\cdot5{,}5=-96{,}20,\\
n_z&=e_{1x}e_{2y}-e_{1y}e_{2x}=\left(-2{,}0\right)\cdot1{,}7-1{,}1\cdot6{,}7=-10{,}77,
\end{align*}
es decir $\vect{n}=\begin{pmatrix} 33{,}2 \\ -96{,}2 \\ -10{,}8 \end{pmatrix}$ de norma $102{,}4$; el vector hacia el centro es $\vect{c}_t=\begin{pmatrix} 868{,}4 \\ 2054{,}2 \\ 6768{,}5 \end{pmatrix}$ ($\lVert\vect{c}_t\rVert=7126{,}4$), y
\begin{equation}
\lvert\hat{\vect{n}}\cdot\hat{\vect{c}}_t\rvert=\frac{241\,636}{(102{,}4)(7126{,}4)}=0{,}331>0{,}087,\qquad \frac{Z_{\max}}{Z_{\min}}=1{,}010<1{,}12,
\end{equation}
así que el triángulo \emph{se conserva}. Con un triángulo del salto entre el sofá y la pared, en los píxeles $(344,370)$, $(342,372)$, $(344,372)$,
$\vect{e}_1=\begin{pmatrix} 438{,}4 \\ -154 \\ -2492{,}4 \end{pmatrix}$ y $\vect{e}_2=\begin{pmatrix} -1{,}7 \\ 18{,}7 \\ 9{,}7 \end{pmatrix}$ y
\begin{align*}
n_x&=e_{1y}e_{2z}-e_{1z}e_{2y}=\left(-154{,}0\right)\cdot9{,}7-\left(-2492{,}4\right)\cdot18{,}7=45\,114{,}1,\\
n_y&=e_{1z}e_{2x}-e_{1x}e_{2z}=\left(-2492{,}4\right)\cdot\left(-1{,}7\right)-438{,}4\cdot9{,}7=-15{,}40,\\
n_z&=e_{1x}e_{2y}-e_{1y}e_{2x}=438{,}4\cdot18{,}7-\left(-154{,}0\right)\cdot\left(-1{,}7\right)=7\,936{,}3,
\end{align*}
es decir $\vect{n}=\begin{pmatrix} 45114{,}1 \\ -15{,}4 \\ 7936{,}3 \end{pmatrix}$ de norma $45806{,}8$; con $\vect{c}_t=\begin{pmatrix} -3248 \\ 1219{,}5 \\ 17930{,}9 \end{pmatrix}$ ($\lVert\vect{c}_t\rVert=18263{,}5$),
\begin{equation}
\lvert\hat{\vect{n}}\cdot\hat{\vect{c}}_t\rvert=\frac{4\,244\,669}{(45806{,}8)(18263{,}5)}=0{,}005<0{,}087,\qquad \frac{Z_{\max}}{Z_{\min}}=1{,}58>1{,}12,
\end{equation}
y el triángulo \emph{se elimina}. De los 305\,602 triángulos se eliminaron 2\,612 (0{,}85\,\%) y quedaron 302\,990 en la capa de adelante.

\begin{figure}[htbp]
\centering
\includegraphics[width=0.97\linewidth]{hab_corte.jpg}
\caption{Izquierda: profundidad a lo largo de la fila $v=370$ de la foto; el salto del sofá a la pared del fondo es donde se cortan los triángulos. Derecha:
$|\hat{\vect{n}}\cdot\hat{\vect{c}}_t|$ de un triángulo del piso (se conserva) y de uno del salto (se elimina) frente al umbral $\operatorname{sen}5^\circ$.}
\label{fig:hab-corte}
\end{figure}

\paragraph{Paso 5: la capa de atrás.} Detrás de los muebles se agrega una segunda capa de triángulos con la profundidad del cuarto vacío, y la
red LaMa inventa su textura. En esta sala la capa de atrás aportó 42\,168 vértices a los 153\,600 de la capa de adelante, y la matriz final tiene $N=195\,768$ vértices ($V\in\mathbb{R}^{195\,768\times3}$) y
$M=380\,493$ triángulos ($F\in\mathbb{N}^{380\,493\times3}$).

\paragraph{Paso 6: enderezar la escena con una rotación.} El celular casi nunca está derecho, y la normal del piso medida no es exactamente vertical. La del piso es
$\hat{\vect{n}}=\begin{pmatrix} 0{,}00963 \\ -0{,}99828 \\ -0{,}05776 \end{pmatrix}$ y la vertical es $\vect{m}=\begin{pmatrix} 0 \\ -1 \\ 0 \end{pmatrix}$ (hacia arriba, en los ejes de la cámara). Se busca la rotación $R$ que lleva $\hat{\vect{n}}$ a $\vect{m}$ con la
fórmula de Rodrigues, $R=I+[\vect{v}]_\times+[\vect{v}]_\times^{2}/(1+c)$, donde
\begin{equation}
\vect{v}=\hat{\vect{n}}\times\vect{m}=\begin{pmatrix} -0{,}05776 \\ 0 \\ -0{,}00963 \end{pmatrix},\qquad c=\hat{\vect{n}}\cdot\vect{m}=0{,}998284,
\end{equation}
y el ángulo del giro es $\theta=\arccos c=3{,}36^\circ$ (es decir, la cámara estaba inclinada $3{,}36^\circ$). Las matrices son
\begin{equation}
[\vect{v}]_\times=\begin{pmatrix} 0 & 0{,}00963 & 0 \\ -0{,}00963 & 0 & 0{,}05776 \\ 0 & -0{,}05776 & 0 \end{pmatrix},\qquad [\vect{v}]_\times^{2}=\begin{pmatrix} -0{,}00009 & 0 & 0{,}00056 \\ 0 & -0{,}00343 & 0 \\ 0{,}00056 & 0 & -0{,}00334 \end{pmatrix} ,
\end{equation}
y, sumando con $1/(1+c)=0{,}5004$,
\begin{equation}
R=\begin{pmatrix} 0{,}99995 & 0{,}00963 & 0{,}00028 \\ -0{,}00963 & 0{,}99828 & 0{,}05776 \\ 0{,}00028 & -0{,}05776 & 0{,}99833 \end{pmatrix} .
\end{equation}
La comprobación es que $R\hat{\vect{n}}=\begin{pmatrix} 0 \\ -1 \\ 0 \end{pmatrix}=\vect{m}$, que $\det R=1{,}000000$ y que $R^{T}R=I$ (coincide con la identidad con diferencias menores que $10^{-8}$): la rotación no
deforma el cuarto. Después se compone con el cambio de ejes del visor $R_o=\operatorname{diag}(1,-1,-1)$ y todos los vértices se giran con un solo producto de matrices,
\begin{equation}
V_{\text{giro}}=V\,G^{T},\qquad G=R_oR=\begin{pmatrix} 0{,}99995 & 0{,}00963 & 0{,}00028 \\ 0{,}00963 & -0{,}99828 & -0{,}05776 \\ -0{,}00028 & 0{,}05776 & -0{,}99833 \end{pmatrix} .
\end{equation}
Los puntos del piso tenían una dispersión de la coordenada vertical de $67{,}6$ mm y quedan con $6{,}5$ mm, alrededor de la altura $y=-1\,207$ mm, que es la altura de la cámara sobre el
piso ($1{,}21$ m, igual a $|d|$). Los vectores $\hat{\vect{n}}$, $\vect{m}$ y $\vect{v}$ y las cifras anteriores están en la figura~\ref{fig:vec-rodrigues}, y el antes y el después del piso, en la figura~\ref{fig:hab-rotacion}.

\begin{figure}[htbp]
\centering
\includegraphics[width=0.97\linewidth]{vec_rodrigues.jpg}
\caption{Los vectores de la rotación de Rodrigues: la normal medida del piso $\hat{\vect{n}}$, la vertical $\vect{m}$ y el eje de giro $\vect{v}=\hat{\vect{n}}\times\vect{m}$, perpendicular a los dos.}
\label{fig:vec-rodrigues}
\end{figure}

\begin{figure}[htbp]
\centering
\includegraphics[width=\linewidth]{hab_rotacion.jpg}
\caption{Puntos del piso antes de girar (inclinados, en los ejes de la cámara), después de girar con $G=R_oR$ (horizontales) y distribución de su altura.}
\label{fig:hab-rotacion}
\end{figure}

\paragraph{Los vectores de la habitación en un cuadro.} El cuadro~\ref{tab:vec-hab} reúne los vectores que aparecieron y la operación con la que se
obtuvo cada uno.

\begin{table}[htbp]
\caption{Vectores del proceso de una habitación y operación con la que se obtienen}
\label{tab:vec-hab}
\centering\small
\begin{tabular}{@{}p{2.5cm}p{4.0cm}p{3.1cm}p{4.2cm}@{}}
\toprule
Vector & Qué representa & Valor en la sala & Operación \\
\midrule
$\vect{r}=K^{-1}\vect{p}$ & dirección del rayo de un píxel & $\begin{pmatrix} 0{,}2004 \\ 0{,}2681 \\ 1 \end{pmatrix}$ (píxel del piso) & producto matriz--vector; norma $1{,}0545$ \\
$\vect{P}=Z\,\vect{r}$ & punto 3D del píxel & $\begin{pmatrix} 0{,}75 \\ 1 \\ 3{,}73 \end{pmatrix}$ m & producto por un escalar \\
$\vect{e}_1,\ \vect{e}_2$ & dos lados del triángulo de RANSAC & $\lVert\vect{e}_1\rVert=3{,}01$ m & resta de puntos \\
$\hat{\vect{n}}$ & normal de un plano & $\begin{pmatrix} -0{,}0295 \\ -0{,}9974 \\ -0{,}0657 \end{pmatrix}$ & producto cruz y norma \\
$\hat{\vect{n}}_i\cdot\hat{\vect{n}}_j$ & ángulo entre dos planos & $0{,}0066$ (piso--pared) & producto punto \\
$t_k=d_k/(\hat{\vect{n}}_k\cdot\vect{r})$ & corte del rayo con un plano & $4{,}85$ m (detrás del sofá) & producto punto y división \\
$\vect{v}=\hat{\vect{n}}\times\vect{m}$ & eje del giro que endereza la sala & $\begin{pmatrix} -0{,}058 \\ 0 \\ -0{,}010 \end{pmatrix}$ & producto cruz \\
$\hat{\vect{n}}\cdot\hat{\vect{c}}_t$ & si un triángulo se ve de canto & $0{,}331$ (piso), $0{,}005$ (salto) & producto punto con el rayo al centro \\
\bottomrule
\end{tabular}
\end{table}
